\section{Criptomonedas}

Las criptomonedas sin un sistema digital de intercambio
descentralizado (peer to peer) en el que la criptografia es usada
para generar y distribuir esta moneda digital.
Al no existir ser descentralizado, no existe autoridad u organizacion
que regule a una criptomoneda, por lo que para las transacciones es
necesaria algun tipo de verificacion durante este proceso.

\begin{itemize}
  \item Que el monto sea correcto
  \item Que el remitente tenga los fondos necerarios para
    la operacion
  \item Asegurarse de que no haya duplicamiento de la unidad (moneda)
\end{itemize}

El proceso de aseguramiento y verificacion durante la transaccion es
parte de lo que se conoce como "mineria".

\subsection{Tipos de criptomonedas}

El top de criptomonedas mas valiosas varia bastante con el tiempo,
asi como la lista de empresas mas valiosas que cotizan en el mercado,
esta tabla enlista aquellas mas representativas.

\begin{table}[ht]
  \centering
  \caption{Comparativa de Criptomonedas Principales}
  \begin{tabularx}{\textwidth}{XXXX} % Sin líneas verticales (|)
    \toprule % Línea gruesa superior
    \textbf{Bitcoin} & \textbf{Ethereum} & \textbf{Ripple} &
    \textbf{Litecoin} \\ \midrule % Línea media

    Primera criptomoneda digital con fecha de publicación del 2009,
    registra las transacciones en el ``Blockchain''. &
    Lanzada en 2015, plataforma que permite contratos inteligentes y
    aplicaciones descentralizadas. &
    Red de liquidación global en tiempo real que ofrece pagos
    internacionales a bajo costo. &
    Basada en código abierto, con una tasa de generación de bloques
    más alta que Bitcoin. \\

    \bottomrule % Línea inferior
  \end{tabularx}
\end{table}

Siendo Bitccoin la moneda mas representativa o valorada en este
mercado, se usara de ejemplo en este documento para dar expolicacion
sobre el fundamento, creacion y mantenimiento de las cxcriptomonedas
de manera general.

\subsection{Blockain}

El blockchain de Bitcoin basicamente es un registro publico quye
aloja todas las transacciones hechas, ademas de incluir la creacion
de nuevas unidades. Se crea una cadena de bloques mediante la
conexion del bloque nuevo con la referencia de los anteriores.

Los componentes principales de un bloque son:

\begin{itemize}
  \item Un valor hash que referencia al bloque inmediato
    que lo precede
  \item Referencia propia
  \item Fecha de creacion
  \item Informacion: Transacciones hechas durante la
    creacion del bloque
\end{itemize}

\subsection{Como se sustenta el sistema}

Para entender el mecanismo que soporta el Bitcoin, es necesario
primero hablar sobre el rol de un minero en este sistema. Un minero
junta o recolecta transacciones de Bitcoin, ayuda a verificar su
legitimidad (monto correcto, fondos suficientes, sin duplicados) y
los almacena como ‘candidatos a bloque’.

Quien quiera puede convertirse en minero, el único requisito es
descargar el software y una copia reciente del Blockchain.
Actualmente la carrera por la nueva creación de bloques se ha vuelto
compleja debido a la alta demanda computacional lo que lleva a la
creación de granjas de minería con hardware altamente especializado,
lo que se traduce en alto consumo eléctrico (para mantener el
hardware prendido y enfriarlo) y una huella enorme de carbón.

Para que un bloque sea aceptado en el Blockchain, es necesario que se
cumplan criterios especificos:

\begin{itemize}
  \item Todas las transacciones registradas tienen que ser LEGITIMAS
  \item Fingerprint (Huella digital valida)
\end{itemize}

Este `fingerprint` se logra al sacar el valor hash del bloque
candidato usando la funcion hash dSHA256, no basta con obtener
cualquier valor aleatorio, se requiere que esta huella tenga una
caracteristica que lo haga unico (dificil de encontrar).

\begin{itemize}
  \item Valor debajo de un limite
  \item Se deben de mostrar multiples ceros al inicio de la cadena
\end{itemize}

Los mineros constantemente se encuentran resolviendo este problema de
buscar entradas que devuelvan estos valores hash de alta rareza. Si
un minero llega a encontrar un valor aceptado, se añade el bloque
inmediatamente al sistema y quienes participan pueden verificar por
sí mismos que el bloque sea válido.

Por cada bloque nuevo añadido al Blockchain, el sistema crea nuevos
Bitcoins, y todo aquel participante recibe en proporción un premio en
forma de una fracción de Bitcoin. a su vez, por cada ciertos bloques
creados, el sistema disminuye gradualmente la cantidad de nuevos
Bitcoin creados, este tiene un tope máximo de 21 millones de monedas

\subsection{Volatilidad del mercado}

La volatilidad extrema de las criptomonedas, en particular Bitcoin,
puede explicarse por la falta de un valor intrínseco o fundamental
que pueda servir como un ancla para su precio. en cambio, los precios
se impulsan por el sentimiento del mercado, la especulación y las
narrativas de los medios, lo que lleva a la formación de burbujas ya
un compartimiento de efecto de rebaño entre los inversores

\section{Optimizacion de Portafolios}

De manera convencional, el manejo de portafolios se refiere al
proceso de identificar, seleccionar y manejar un conjunto de
inversiones. Estos procesos implican la selección de activos,
asignación de recursos, seguimiento del rendimiento a través del
tiempo y la actualización del portafolio.

Los orígenes de este problema de optimización de portafolios hecha
con un modelo matemático, datan desde las contribuciones que Harry
Markowitz hizo en el año 1952 con su gran aporte que fue parteaguas
para la economía financiera, formando la fundación de lo que hoy se
conoce como Teoría del Portafolio Moderno. Desde entonces, este campo
se ha vuelto inmenso.

\begin{center}
  \itshape
  ``Optimización de portafolio, se puede definir como el uso de
  métodos matemáticos que apoya la selección de portafolios de
  activos a elección teniendo en cuenta la preferencia de toma de
  decisiones, límites relevantes e incertidumbre''
\end{center}

Este problema matemático puede ser formulado de distintas maneras,
pero los problemas principales se resumen en:

\begin{enumerate}[label=\Roman*., leftmargin=1.5cm]
  \item Minimizar el riesgo para un rendimiento establecido
  \item Maximizar el rendimiento esperado dado un riesgo establecido
  \item Minimizar el riesgo y maximizar el rendimiento esperado
    usando un factor de rechazo de riesgo especificado
  \item Minimizar el riesgo sin importar el rendimiento esperado
  \item Maximizar el rendimiento sin importar el riesgo
\end{enumerate}

\subsection{Definiciones}
\subsubsection{Rendimento del activo:}
Expresado como la tasa de rendimiento que se define durante un
periodo de tiempo. Una perdida es representada como un valor negativo

\subsubsection{Portafolio:}
Una colección de dos o más activos. El rendimiento esperado (por
periodo) se define
como la suma ponderada del rendimiento de cada activo, sujeto a que
la suma total sea 1:
\begin{gather}
  R_p = \sum_{i=1}^{n} w_i R_i \\
  \sum_{i=1}^{n} w_i = 1
\end{gather}

\subsubsection{Short selling:}
Es un proceso por el cual un inversionista que no posee un activo,
define una posición en el mercado vendiendo el activo por anticipado,
lo que provoca en una caída del precio. Esto matemáticamente se
representa tomando el valor del activo como negativo. Para el modelo
de Markowitz se excluye esta práctica.

\subsubsection{Portafolio eficiente:}
Un portafolio factible es eficiente si tiene el máximo rendimiento
esperado entre todos los portafolios con la misma varianza o
alternativamente si tiene el menor riesgo entre todos los portafolios
que tienen el mismo rendimiento esperado

\subsubsection{Frontera eficiente:}
Es un concepto fundamental que permite visualizar las contras y
ventajas entre riesgo/rendimiento, representa el conjunto de
portafolios que proveen el rendimiento mas alto de acuerdo a un
riesgo dado, o alternativamente llos portafolios que tienen un riesgo
minimo segun un rendimiento dado.

\subsection{Indicadores de riesgo}

\subsubsection{Varianza:}
Dispersión de los rendimientos; mientras mayor
sea, más inciertos son los resultados.
\subsubsection{Semivarianza:}Solo mide la dispersión hacia abajo (riesgo
de pérdidas).
\subsubsection{Desviacion estandar:}Mide la volatilidad total de los activos.
\subsubsection{Value at Risk (VaR):}Pérdida máxima esperada dado un cierto
nivel de confianza.
\subsubsection{Conditional Value at Risk (CVaR):} Pérdida promedio en los
peores casos.

\subsection{Formas de representar un portafolio}
Usar el peso por cada activo no trabajan bien, esto se debe a que con
esta representación, frecuentemente se producen soluciones no
factibles en el cual la suma total del portafolio (Inversión hecha
con el presupuesto dado) no es 1, y esto es una restricción básica de
este modelo

La necesidad de representar de distintas maneras un portafolio se
debe a que las restricciones que se consideren complejas se puedan
incluir de manera mas natural, ademas de que se puede permitir
explorar un espacio de búsqueda aún más amplio, y también se pueden
adaptar diferentes objetivos que van más allá del rendimiento-riesgo

\subsubsection{Campo de bits para cada activo}
El primer bit de la cadena tiene un propósito especial: indica si la
posición en este activo es larga o corta, lo que permite al inversor
obtener ganancias tanto si el precio sube como si baja. Los k bits
restantes no representan el precio del activo directamente, sino que
funcionan como un índice que apunta a un valor específico en una
tabla, que contiene los posibles pesos que el activo puede tener en
el portafolio

\subsubsection{Tipo mochila (Knapsack)}
Está inspirada en el famoso problema de la mochila, Aquí se utiliza
un bit para cada activo que se considera: si el bit es 1, el activo
se incluye en el portafolio y si es 0, se excluye. Esto simplifica el
problema a enfocarse solo en la selección de activos, sin considerar
sus pesos exactos

\subsubsection{Codigos gray}
Variación de la codificación binaria que asegura que solo un bit
cambia entre dos valores numéricos consecutivos,. Puede ayudar a
mejorar la eficiencia de los algoritmo evolutivos el evitar saltos
grandes en el espacio de búsqueda

\subsubsection{Representacion hibrida}
Combina dos partes distintas de un cromosoma:

\begin{itemize}
  \item Primer parte contiene una secuencia que determina
    las identidades de los activos que se incluyen en el portafolio.
    Es decir esta parte decide que activos se eligen
  \item Después de que la primer parte ha
    seleccionado los activos, esta segunda determina los pesos o la
    proporción de la inversión que se asignará a cada uno de ellos
\end{itemize}

\subsection{Reestricciones del mundo real}
Inicialmente el modelo de Markowitz consideraba sólo una restricción,
la suma del peso de los activos debe de ser 1, en otras palabras, se
tiene que usar todo el capital.

\subsubsection{Limites minimos y maximos}
Se imponen límites inferiores y/o superiores en los valores de cada
peso activo, esto significa que un activo no puede representar menos
o más que alguna proporción del capital total
\begin{itemize}
  \item Limite Inferior: No dedicar porcentajes muy pequeños del
    capital a muchos valores (debido a que se generaron altos costos
    de transacción)
  \item Limite Superior: No asignar una proporción grande del capital
    total a un activo, con el fin de minimizar el riesgo (Obligando a
    diversificar)
\end{itemize}

\subsubsection{Cardinalidad}
Esta restricción obliga al número de activos seleccionados en un
portafolio. Existen dos versiones:
\begin{itemize}
  \item Exacta: Impone que el número de valores seleccionados sea
    igual a un valor K
  \item Suave: Proporciona limites inferiores e superiores para el
    número de activos seleccionados, por ejemplo, al menos 5 y no más de 10
\end{itemize}

Cuando se impone esta restricción, no hay un método exacto para
resolver el problema de optimización de portafolios, la frontera de
pareto puede ser discontinua y/o no convexa, lo que causa
dificultades si se usan algún acercamiento basado en funciones lineales

Las estrategias más comunes para manejar la cardinalidad es usar
alguna función que penalice, también de los portafolios con una
cardinalidad alta, se eliminan los activos con menor peso hasta
cumplir con este valor K

\subsubsection{Rotacion}
Se refiere al cambio en el peso de los activos, respecto a una
asignación anterior. Es útil cuando se considera una inversión
durante un periodo largo. En el modelo tradicional, sin restricciones
adicionales, el algoritmo puede proponer grandes cambios en la
composición de los portafolios, y esto puede ser costoso porque se
generan gastos  de comisión o impuestos.

El objetivo es reducir costos por transacciones, mantener estabilidad
en la estrategia y evitar cambios bruscos. La rotación se calcula
como la suma de las diferencias absolutas entre los pesos nuevos y
los anteriores. Si esta suma excede de un límite predefinido, la
solución no es aceptada

\section{Metaheuristica}
La metaheurística es un campo de la informática y la optimización que
se enfoca en el desarrollo de algoritmos de alto nivel para
encontrar, de manera eficiente, soluciones aproximadas a problemas de
optimización complejos, especialmente aquellos para los que no se
conoce un algoritmo exacto o el tiempo de ejecución de un algoritmo
exacto es largo. A diferencia de los métodos heurísticos simples, una
metaheurística está diseñada para guiar el proceso de búsqueda a
través del espacio de soluciones, evitando quedarse atrapada en
óptimos locales y permitiendo una variedad más amplia de soluciones posibles.

\subsection{Algoritmos bioinspirados}
Toman su inspiración de procesos y comportamientos observados en la
naturaleza y la biología, tales como el proceso de evolución, el
comportamiento de enjambre, forrajeo (actividad de buscar alimentos)
y el crecimiento de las plantas La idea principal detrás de estos
algoritmos es imitar la eficiencia de la naturaleza para resolver
problemas complejos.

Tras un rápido crecimiento tecnológico, la complejidad computacional
ha incrementado para resolver problemas del mundo real como:
optimización,  clasificación, planeación y control. Estos
problemas se caracterizan por su alta dimensionalidad, no linealidad,
entornos dinámicos, los métodos tradicionales de optimización suelen
quedar cortos contra estos problemas, debido a ser rígidos, las
limitantes particulares son las combinaciones a gran escala o
espacios de búsqueda enormes, donde la adaptabilidad y exploración
son críticos.

\section{Optimizacion por enjambre de particulas (PSO)}
Optimización de Enjambre de Partículas PSO
Desarrollada por Kennedy y Eberhart, esta metaheurística tiene
inspiración del comportamiento social observado en grupos de
individuos dentro de: parvadas o cardúmenes. Este comportamiento se
basa en la transmisión del suceso de cada individuo a los demás del
grupo, esto resulta en satisfacer más rápido sus necesidades, como la
localización de alimento o un lugar de cobijo.

El algoritmo de PSO identifica la solución óptima explorando
exhaustivamente el espacio de la función objetivo a través del ajuste
sistemático de las rutas de cada individuo, llamado partícula. Cada
partícula en el enjambre tiene propiedades inherentes como la
posición y la velocidad, que les ayudan a moverse hacia una posición
óptima, esta partícula tiene una tendencia a moverse aleatoriamente
con la probabilidad de visitar un lugar inexplorado. Pero al mismo
tiempo, cada partícula es atraída hacia su mejor posición personal
(hasta la iteración actual), así como la mejor posición global del
enjambre. El objetivo de PSO es encontrar la posición óptima en el
espacio de búsqueda, considerando las mejores soluciones asociadas a
cada partícula, este proceso termina cuando no hay mejora del valor
de la función objetivo o se alcanza un límite de iteraciones.

\subsection{Pseudocódigo}

\subsection{Modelos de velocidad}
\subsection{Topologias de comunicación}
